% Shared preamble for learn-attn documents
% Usage: % Shared preamble for learn-attn documents
% Usage: % Shared preamble for learn-attn documents
% Usage: \input{preamble} at the top of each document

\documentclass[11pt, a4paper]{article}

\usepackage[T1]{fontenc}
\usepackage{lmodern}
\usepackage{microtype}
\usepackage[margin=1in]{geometry}
\usepackage{amsmath, amssymb}
\usepackage{booktabs}
\usepackage{enumitem}
\usepackage{parskip}
\usepackage[dvipsnames]{xcolor}
\usepackage{listings}
\usepackage{courier}
\usepackage{pgfplots}
\pgfplotsset{compat=1.18}
\usetikzlibrary{backgrounds}
\usepackage[colorlinks=true, linkcolor=NavyBlue, urlcolor=NavyBlue, citecolor=NavyBlue]{hyperref}

% Code listing style
\definecolor{codebg}{gray}{0.96}
\definecolor{codeframe}{gray}{0.8}
\definecolor{codecomment}{RGB}{90,130,90}
\definecolor{codestring}{RGB}{160,60,60}
\definecolor{codekw}{RGB}{40,40,120}

\lstset{
  language=Python,
  basicstyle=\ttfamily\small,
  backgroundcolor=\color{codebg},
  frame=single,
  rulecolor=\color{codeframe},
  framesep=6pt,
  xleftmargin=6pt,
  xrightmargin=6pt,
  breaklines=true,
  breakatwhitespace=true,
  showstringspaces=false,
  keywordstyle=\color{codekw}\bfseries,
  commentstyle=\color{codecomment}\itshape,
  stringstyle=\color{codestring},
  numbers=none,
  tabsize=4,
  aboveskip=1em,
  belowskip=1em,
  morekeywords={self, None, True, False},
}

% Inline code
\newcommand{\code}[1]{\texttt{#1}}

% Document series header
\newcommand{\docheader}[2]{%
  \begin{center}
    {\Large\bfseries #1}\\[6pt]
    {\itshape Document #2 of 6\,---\,Building a GPT from Scratch}
  \end{center}
  \vspace{1em}
}

% --- Code block annotations ---
% Small colored tags placed before a listing to clarify its role.

% Code that the reader should put into a project file.
% Usage: \filetag{src/model.py}  (before \begin{lstlisting})
\newcommand{\filetag}[1]{%
  \par\noindent
  \colorbox{black!8}{\small\sffamily\bfseries File: \code{#1}}%
  \nopagebreak\par\nopagebreak\vspace{-0.6em}%
}

% Standalone demo — try in a Python REPL or scratch script, not part of the project.
\newcommand{\demotag}{%
  \par\noindent
  \colorbox{NavyBlue!12}{\small\sffamily\bfseries Demo \textnormal{--- try in a REPL}}%
  \nopagebreak\par\nopagebreak\vspace{-0.6em}%
}

% Pseudocode or conceptual snippet illustrating an idea.
\newcommand{\concepttag}{%
  \par\noindent
  \colorbox{black!6}{\small\sffamily\bfseries Conceptual \textnormal{--- illustrating the idea}}%
  \nopagebreak\par\nopagebreak\vspace{-0.6em}%
}

% Expected terminal or REPL output.
\newcommand{\outputtag}{%
  \par\noindent
  \colorbox{black!6}{\small\sffamily\bfseries Expected output}%
  \nopagebreak\par\nopagebreak\vspace{-0.6em}%
}

% Shell command the reader should run.
% Usage: \shelltag  (before a listing with language=bash or similar)
\newcommand{\shelltag}{%
  \par\noindent
  \colorbox{Goldenrod!20}{\small\sffamily\bfseries Shell \textnormal{--- run in your terminal}}%
  \nopagebreak\par\nopagebreak\vspace{-0.6em}%
}

% Verification checkpoint — run this and compare.
\newcommand{\checkpoint}[1]{%
  \medskip\par\noindent
  \fbox{\parbox{\dimexpr\linewidth-2\fboxsep-2\fboxrule}{%
    {\sffamily\bfseries Checkpoint:} #1}}%
  \medskip
}

% Tip — clarifies a common point of confusion.
\newcommand{\tip}[1]{%
  \medskip\par\noindent
  \colorbox{ForestGreen!10}{\parbox{\dimexpr\linewidth-2\fboxsep}{%
    {\sffamily\bfseries\color{ForestGreen!70!black} Tip:} #1}}%
  \medskip
}
 at the top of each document

\documentclass[11pt, a4paper]{article}

\usepackage[T1]{fontenc}
\usepackage{lmodern}
\usepackage{microtype}
\usepackage[margin=1in]{geometry}
\usepackage{amsmath, amssymb}
\usepackage{booktabs}
\usepackage{enumitem}
\usepackage{parskip}
\usepackage[dvipsnames]{xcolor}
\usepackage{listings}
\usepackage{courier}
\usepackage{pgfplots}
\pgfplotsset{compat=1.18}
\usetikzlibrary{backgrounds}
\usepackage[colorlinks=true, linkcolor=NavyBlue, urlcolor=NavyBlue, citecolor=NavyBlue]{hyperref}

% Code listing style
\definecolor{codebg}{gray}{0.96}
\definecolor{codeframe}{gray}{0.8}
\definecolor{codecomment}{RGB}{90,130,90}
\definecolor{codestring}{RGB}{160,60,60}
\definecolor{codekw}{RGB}{40,40,120}

\lstset{
  language=Python,
  basicstyle=\ttfamily\small,
  backgroundcolor=\color{codebg},
  frame=single,
  rulecolor=\color{codeframe},
  framesep=6pt,
  xleftmargin=6pt,
  xrightmargin=6pt,
  breaklines=true,
  breakatwhitespace=true,
  showstringspaces=false,
  keywordstyle=\color{codekw}\bfseries,
  commentstyle=\color{codecomment}\itshape,
  stringstyle=\color{codestring},
  numbers=none,
  tabsize=4,
  aboveskip=1em,
  belowskip=1em,
  morekeywords={self, None, True, False},
}

% Inline code
\newcommand{\code}[1]{\texttt{#1}}

% Document series header
\newcommand{\docheader}[2]{%
  \begin{center}
    {\Large\bfseries #1}\\[6pt]
    {\itshape Document #2 of 6\,---\,Building a GPT from Scratch}
  \end{center}
  \vspace{1em}
}

% --- Code block annotations ---
% Small colored tags placed before a listing to clarify its role.

% Code that the reader should put into a project file.
% Usage: \filetag{src/model.py}  (before \begin{lstlisting})
\newcommand{\filetag}[1]{%
  \par\noindent
  \colorbox{black!8}{\small\sffamily\bfseries File: \code{#1}}%
  \nopagebreak\par\nopagebreak\vspace{-0.6em}%
}

% Standalone demo — try in a Python REPL or scratch script, not part of the project.
\newcommand{\demotag}{%
  \par\noindent
  \colorbox{NavyBlue!12}{\small\sffamily\bfseries Demo \textnormal{--- try in a REPL}}%
  \nopagebreak\par\nopagebreak\vspace{-0.6em}%
}

% Pseudocode or conceptual snippet illustrating an idea.
\newcommand{\concepttag}{%
  \par\noindent
  \colorbox{black!6}{\small\sffamily\bfseries Conceptual \textnormal{--- illustrating the idea}}%
  \nopagebreak\par\nopagebreak\vspace{-0.6em}%
}

% Expected terminal or REPL output.
\newcommand{\outputtag}{%
  \par\noindent
  \colorbox{black!6}{\small\sffamily\bfseries Expected output}%
  \nopagebreak\par\nopagebreak\vspace{-0.6em}%
}

% Shell command the reader should run.
% Usage: \shelltag  (before a listing with language=bash or similar)
\newcommand{\shelltag}{%
  \par\noindent
  \colorbox{Goldenrod!20}{\small\sffamily\bfseries Shell \textnormal{--- run in your terminal}}%
  \nopagebreak\par\nopagebreak\vspace{-0.6em}%
}

% Verification checkpoint — run this and compare.
\newcommand{\checkpoint}[1]{%
  \medskip\par\noindent
  \fbox{\parbox{\dimexpr\linewidth-2\fboxsep-2\fboxrule}{%
    {\sffamily\bfseries Checkpoint:} #1}}%
  \medskip
}

% Tip — clarifies a common point of confusion.
\newcommand{\tip}[1]{%
  \medskip\par\noindent
  \colorbox{ForestGreen!10}{\parbox{\dimexpr\linewidth-2\fboxsep}{%
    {\sffamily\bfseries\color{ForestGreen!70!black} Tip:} #1}}%
  \medskip
}
 at the top of each document

\documentclass[11pt, a4paper]{article}

\usepackage[T1]{fontenc}
\usepackage{lmodern}
\usepackage{microtype}
\usepackage[margin=1in]{geometry}
\usepackage{amsmath, amssymb}
\usepackage{booktabs}
\usepackage{enumitem}
\usepackage{parskip}
\usepackage[dvipsnames]{xcolor}
\usepackage{listings}
\usepackage{courier}
\usepackage{pgfplots}
\pgfplotsset{compat=1.18}
\usetikzlibrary{backgrounds}
\usepackage[colorlinks=true, linkcolor=NavyBlue, urlcolor=NavyBlue, citecolor=NavyBlue]{hyperref}

% Code listing style
\definecolor{codebg}{gray}{0.96}
\definecolor{codeframe}{gray}{0.8}
\definecolor{codecomment}{RGB}{90,130,90}
\definecolor{codestring}{RGB}{160,60,60}
\definecolor{codekw}{RGB}{40,40,120}

\lstset{
  language=Python,
  basicstyle=\ttfamily\small,
  backgroundcolor=\color{codebg},
  frame=single,
  rulecolor=\color{codeframe},
  framesep=6pt,
  xleftmargin=6pt,
  xrightmargin=6pt,
  breaklines=true,
  breakatwhitespace=true,
  showstringspaces=false,
  keywordstyle=\color{codekw}\bfseries,
  commentstyle=\color{codecomment}\itshape,
  stringstyle=\color{codestring},
  numbers=none,
  tabsize=4,
  aboveskip=1em,
  belowskip=1em,
  morekeywords={self, None, True, False},
}

% Inline code
\newcommand{\code}[1]{\texttt{#1}}

% Document series header
\newcommand{\docheader}[2]{%
  \begin{center}
    {\Large\bfseries #1}\\[6pt]
    {\itshape Document #2 of 6\,---\,Building a GPT from Scratch}
  \end{center}
  \vspace{1em}
}

% --- Code block annotations ---
% Small colored tags placed before a listing to clarify its role.

% Code that the reader should put into a project file.
% Usage: \filetag{src/model.py}  (before \begin{lstlisting})
\newcommand{\filetag}[1]{%
  \par\noindent
  \colorbox{black!8}{\small\sffamily\bfseries File: \code{#1}}%
  \nopagebreak\par\nopagebreak\vspace{-0.6em}%
}

% Standalone demo — try in a Python REPL or scratch script, not part of the project.
\newcommand{\demotag}{%
  \par\noindent
  \colorbox{NavyBlue!12}{\small\sffamily\bfseries Demo \textnormal{--- try in a REPL}}%
  \nopagebreak\par\nopagebreak\vspace{-0.6em}%
}

% Pseudocode or conceptual snippet illustrating an idea.
\newcommand{\concepttag}{%
  \par\noindent
  \colorbox{black!6}{\small\sffamily\bfseries Conceptual \textnormal{--- illustrating the idea}}%
  \nopagebreak\par\nopagebreak\vspace{-0.6em}%
}

% Expected terminal or REPL output.
\newcommand{\outputtag}{%
  \par\noindent
  \colorbox{black!6}{\small\sffamily\bfseries Expected output}%
  \nopagebreak\par\nopagebreak\vspace{-0.6em}%
}

% Shell command the reader should run.
% Usage: \shelltag  (before a listing with language=bash or similar)
\newcommand{\shelltag}{%
  \par\noindent
  \colorbox{Goldenrod!20}{\small\sffamily\bfseries Shell \textnormal{--- run in your terminal}}%
  \nopagebreak\par\nopagebreak\vspace{-0.6em}%
}

% Verification checkpoint — run this and compare.
\newcommand{\checkpoint}[1]{%
  \medskip\par\noindent
  \fbox{\parbox{\dimexpr\linewidth-2\fboxsep-2\fboxrule}{%
    {\sffamily\bfseries Checkpoint:} #1}}%
  \medskip
}

% Tip — clarifies a common point of confusion.
\newcommand{\tip}[1]{%
  \medskip\par\noindent
  \colorbox{ForestGreen!10}{\parbox{\dimexpr\linewidth-2\fboxsep}{%
    {\sffamily\bfseries\color{ForestGreen!70!black} Tip:} #1}}%
  \medskip
}

\begin{document}
\docheader{Attention from First Principles}{1}

\section{Project Setup}

Before we begin, create the project and install PyTorch:

\shelltag
\begin{lstlisting}[language=bash]
mkdir learn-attn && cd learn-attn
uv init
uv add torch --index https://download.pytorch.org/whl/cu126
\end{lstlisting}

This gives you a Python environment with PyTorch and CUDA support. All demo code in Documents 01--03 can be run in a Python REPL (\code{uv run python}). Starting in Document 04, we will create the actual project source files.

\checkpoint{Run \code{uv run python -c "import torch; print(torch.cuda.is\_available())"} --- it should print \code{True}.}

This document derives scaled dot-product attention from the ground up, builds multi-head attention on top of it, and introduces the causal mask needed for language modeling. By the end, every line of \code{CausalSelfAttention} in \code{src/model.py} should make sense.

\bigskip\hrule\bigskip

\section{What Problem Does Attention Solve?}

Consider a sequence of $T$ token embeddings, each a vector in $\mathbb{R}^d$. (How tokens become vectors is covered in Document~03, Section~5. For now, just accept that each token is represented as a $d$-dimensional vector.) Every position in the sequence carries some meaning, but meaning in language is deeply contextual --- the word ``bank'' means something different depending on whether ``river'' or ``account'' appears nearby. So each position needs a mechanism to gather information from other positions.

Before attention, the dominant approach was recurrent neural networks (RNNs). An RNN processes a sequence left to right, maintaining a hidden state that gets updated at each step. This has two problems:

\begin{enumerate}
  \item \textbf{Sequential bottleneck.} You cannot compute the output at position $t$ until you have computed positions 1 through $t-1$. This kills parallelism on GPUs.
  \item \textbf{Information decay.} Information from position 1 must survive through every hidden-state update to reach position $T$. In practice, gradients vanish or explode, and the model struggles with long-range dependencies despite techniques like gating (LSTM, GRU).
\end{enumerate}

Attention solves both problems with a single idea: \textbf{let each position compute a weighted average of all positions, where the weights are determined dynamically by the content at those positions.} There is no sequential dependency --- every position can attend to every other position in parallel. And there is no distance decay --- position 1 is just as accessible as position $T-1$, if the content says it should be.

The rest of this document is about making that idea precise.

\bigskip\hrule\bigskip

\section{Dot Products as Similarity}

The mechanism that determines ``how much should position $i$ attend to position $j$?'' is the dot product.

Recall from linear algebra: for two vectors $\mathbf{a}, \mathbf{b} \in \mathbb{R}^d$,

\[
\mathbf{a} \cdot \mathbf{b} = \|\mathbf{a}\| \, \|\mathbf{b}\| \cos(\theta)
\]

where $\theta$ is the angle between them. When $\mathbf{a}$ and $\mathbf{b}$ point in similar directions (small $\theta$), $\cos(\theta)$ is close to 1 and the dot product is large. When they are orthogonal, the dot product is zero. When they point in opposite directions, it is negative.

In attention, we assign each position two roles:

\begin{itemize}
  \item A \textbf{query} vector $\mathbf{q}$ --- ``what am I looking for?''
  \item A \textbf{key} vector $\mathbf{k}$ --- ``what do I contain?''
\end{itemize}

The dot product $\mathbf{q} \cdot \mathbf{k}$ measures how well the query at one position matches the key at another. High dot product means high relevance.

In practice, $\mathbf{q}$ and $\mathbf{k}$ are not stored directly --- they are produced by multiplying the token embedding $\mathbf{x}$ by learned weight matrices: $\mathbf{q} = \mathbf{x}\, W^Q$ and $\mathbf{k} = \mathbf{x}\, W^K$. The network learns these projections during training.

\textbf{Concrete example.} Suppose we have 3 positions with query and key vectors in $\mathbb{R}^4$:

\demotag
\begin{lstlisting}
import torch

Q = torch.tensor([
    [1.0,  0.0,  1.0,  0.0],  # query at position 0
    [0.0,  1.0,  0.0,  1.0],  # query at position 1
    [1.0,  1.0,  0.0,  0.0],  # query at position 2
])

K = torch.tensor([
    [1.0,  0.0,  1.0,  0.0],  # key at position 0
    [0.0,  1.0,  0.0,  1.0],  # key at position 1
    [1.0,  0.0,  0.0,  1.0],  # key at position 2
])

scores = torch.matmul(Q, K.transpose(-2, -1))
print(scores)
\end{lstlisting}

Output:

\outputtag
\begin{lstlisting}
tensor([[2., 0., 1.],
        [0., 2., 1.],
        [1., 1., 1.]])
\end{lstlisting}

Reading row 0: position 0's query is most similar to position 0's key (score 2), completely orthogonal to position 1's key (score 0), and moderately similar to position 2's key (score 1). This is exactly what you would verify by hand --- $\mathbf{q}_0 = [1, 0, 1, 0]$ and $\mathbf{k}_0 = [1, 0, 1, 0]$ share all their nonzero components, while $\mathbf{k}_1 = [0, 1, 0, 1]$ has nonzero components exactly where $\mathbf{q}_0$ is zero.

\tip{Why the transpose? A single dot product $\mathbf{q}_i \cdot \mathbf{k}_j$ gives one similarity score. We need \emph{all} $T^2$ pairwise scores at once. \code{K.transpose(-2, -1)} flips $K$ from $(T, d_k)$ to $(d_k, T)$, turning each key into a column. Then \code{torch.matmul(Q, K\^{}T)} computes row $\times$ column for every $(i, j)$ pair --- entry $(i, j)$ of the result is exactly $\mathbf{q}_i \cdot \mathbf{k}_j$. The matrix multiply is just all $T^2$ dot products in one shot.}

\textbf{Dimensions.} We write $d_k$ for the dimension of each query and key vector (the same $d$ from Section~2). If $Q$ is $(T, d_k)$ and $K$ is $(T, d_k)$, then $QK^\top$ is $(T, d_k) \times (d_k, T) = (T, T)$. Each entry $(i, j)$ in this matrix is the raw attention score from position $i$ to position $j$. We get a score for every pair of positions.

\bigskip\hrule\bigskip

\section{Why Scale by $\sqrt{d_k}$}

Those raw scores have a problem: their magnitude depends on $d_k$, and this matters a lot when we feed them into softmax.

Here is the variance argument. Suppose each component of $\mathbf{q}$ and $\mathbf{k}$ is drawn independently from $\mathcal{N}(0, 1)$. Then:

\begin{itemize}
  \item Each product $q_i k_i$ has mean $\mathbb{E}[q_i k_i] = \mathbb{E}[q_i]\mathbb{E}[k_i] = 0$ (since they are independent and zero-mean).
  \item Each product $q_i k_i$ has variance $\mathrm{Var}(q_i k_i) = \mathbb{E}[q_i^2]\mathbb{E}[k_i^2] = 1 \cdot 1 = 1$.
  \item The dot product $\mathbf{q} \cdot \mathbf{k} = \sum_{i=1}^{d_k} q_i k_i$. Since they are independent, $\mathrm{Var}(\mathbf{q} \cdot \mathbf{k}) = d_k$.
  \item The standard deviation is $\sqrt{d_k}$.
\end{itemize}

So the dot product is a zero-mean random variable with standard deviation $\sqrt{d_k}$. When $d_k = 64$ (a typical value in practice), the standard deviation is 8. Dot products routinely land in the range $[-16, +16]$ or wider.

Why does this matter? Consider what softmax does with large inputs. If one score is 10 and the rest are near 0, then $\exp(10) \approx 22026$ dominates the denominator, and the output is nearly one-hot. In the backward pass, the gradients of softmax near a one-hot output are very close to zero --- the model is ``confident'' and stops learning. This is the softmax saturation problem.

The fix is to divide by $\sqrt{d_k}$ before applying softmax, normalizing the scores back to unit variance:

\demotag
\begin{lstlisting}
import math

d_k = Q.size(-1)
scores = torch.matmul(Q, K.transpose(-2, -1)) / math.sqrt(d_k)
\end{lstlisting}

There is a nice way to think about this: dividing by $\sqrt{d_k}$ is equivalent to setting the softmax \textbf{temperature} to $\sqrt{d_k}$. Higher temperature makes the distribution more uniform (softer attention); lower temperature makes it more peaked. By scaling, we ensure the temperature is calibrated regardless of dimension, so the model starts training with reasonable gradient flow.

\bigskip\hrule\bigskip

\section{Softmax Turns Scores into Weights}

After scaling, we have a $(T, T)$ matrix of real-valued scores. We need to convert each row into a probability distribution --- non-negative values that sum to 1 --- so we can take a weighted average. Softmax does this:

\[
\mathrm{softmax}(z_i) = \frac{\exp(z_i)}{\sum_j \exp(z_j)}
\]

Properties:
\begin{itemize}
  \item Every output is in $(0, 1)$ --- strictly positive, never exactly zero.
  \item The outputs sum to 1 across the dimension we apply softmax over.
  \item It is monotonic: larger inputs get larger probabilities.
  \item It is differentiable everywhere, which is necessary for backpropagation.
\end{itemize}

In code:

\demotag
\begin{lstlisting}
import torch.nn.functional as F

weights = F.softmax(scores, dim=-1)  # (T, T)
\end{lstlisting}

The \code{dim=-1} means we apply softmax across the last dimension (the key dimension), so each row independently becomes a probability distribution. Row $i$ tells us: ``how much should position $i$ attend to each other position?''

\tip{Think of softmax as two steps: (1)~pass every element through $e^x$, which makes all values strictly positive \emph{and} amplifies larger values exponentially; (2)~divide each result by the sum of all of them, so the vector sums to~1. Step~1 is the nonlinearity that separates winners from losers; step~2 is just linear normalization.}

\textbf{Connection to temperature.} More generally, softmax with temperature $\tau$ is $\mathrm{softmax}(\mathbf{z} / \tau)$. When $\tau$ is small, the distribution sharpens toward a one-hot on the maximum element. When $\tau$ is large, the distribution flattens toward uniform. The $\sqrt{d_k}$ scaling is exactly this: we implicitly set $\tau = \sqrt{d_k}$, preventing the distribution from being too sharp at initialization.

\bigskip\hrule\bigskip

\section{The Full Attention Formula}

We now have all the ingredients. Each position also has a \textbf{value} vector $\mathbf{v}$ --- ``what I contribute to others if they attend to me.'' Like queries and keys, it is a learned projection of the token embedding: $\mathbf{v} = \mathbf{x}\, W^V$. The full scaled dot-product attention formula is:

\[
\mathrm{Attention}(Q, K, V) = \mathrm{softmax}\!\left(\frac{QK^\top}{\sqrt{d_k}}\right) V
\]

\textbf{Dimension walkthrough}, step by step:

\begin{tabular}{lll}
\toprule
Expression & Shape & What it represents \\
\midrule
$Q$ & $(T, d_k)$ & Query vectors for each position \\
$K$ & $(T, d_k)$ & Key vectors for each position \\
$V$ & $(T, d_k)$ & Value vectors for each position \\
$QK^\top$ & $(T, T)$ & Raw attention scores for all position pairs \\
$QK^\top / \sqrt{d_k}$ & $(T, T)$ & Scaled scores (unit variance) \\
$\mathrm{softmax}(QK^\top / \sqrt{d_k})$ & $(T, T)$ & Attention weights (each row sums to 1) \\
$\mathrm{softmax}(\cdots)\, V$ & $(T, d_k)$ & Weighted combination of value vectors \\
\bottomrule
\end{tabular}

The output has shape $(T, d_k)$. Each row $i$ of the output is a weighted average of the value vectors, where the weights come from how well position $i$'s query matches each position's key.

\textbf{Complete implementation:}

This standalone function shows the algorithm. Our project uses PyTorch's fused \code{F.scaled\_dot\_product\_attention} instead.

\demotag
\begin{lstlisting}
import math
import torch
import torch.nn.functional as F

def attention(Q, K, V, mask=None):
    """Scaled dot-product attention.

    Args:
        Q: (*, T, d_k) query vectors
        K: (*, S, d_k) key vectors
        V: (*, S, d_k) value vectors
        mask: broadcastable to (*, T, S), 0 where attention is forbidden

    Returns:
        (*, T, d_k) attended values
    """
    d_k = Q.size(-1)
    scores = torch.matmul(Q, K.transpose(-2, -1)) / math.sqrt(d_k)
    if mask is not None:
        scores = scores.masked_fill(mask == 0, float('-inf'))
    weights = F.softmax(scores, dim=-1)
    return torch.matmul(weights, V)
\end{lstlisting}

A few things to note:
\begin{itemize}
  \item The \code{*} in the shape annotations means this works with arbitrary leading batch dimensions --- useful when processing multiple sequences at once.
  \item The \code{mask} argument is for causal masking (Section~8). Where the mask is 0, we set the score to $-\infty$ before softmax, which makes the weight exactly 0.
  \item The function is pure matrix multiplications and elementwise operations --- extremely GPU-friendly.
\end{itemize}

Let us run it. We reuse $Q$ and $K$ from Section~3 and define a simple $V$ where each position's value is a distinct 2D vector:

\demotag
\begin{lstlisting}
V = torch.tensor([
    [1.0, 0.0],  # value at position 0
    [0.0, 1.0],  # value at position 1
    [0.5, 0.5],  # value at position 2
])

out = attention(Q, K, V)
print(out)
\end{lstlisting}

\outputtag
\begin{lstlisting}
tensor([[0.6601, 0.3399],
        [0.3399, 0.6601],
        [0.5000, 0.5000]])
\end{lstlisting}

Interpreting each row:

\begin{itemize}
  \item \textbf{Position 0} --- its query best matches its own key (Section~3), so it pulls mostly from $V_0 = [1, 0]$. The output $[0.66, 0.34]$ leans heavily toward $V_0$.
  \item \textbf{Position 1} --- symmetric case. Its query best matches $K_1$, so it pulls mostly from $V_1 = [0, 1]$. The output $[0.34, 0.66]$ leans toward $V_1$.
  \item \textbf{Position 2} --- its query matches all three keys roughly equally (scores were all~1 in Section~3), so attention is nearly uniform. The output $[0.50, 0.50]$ is close to the simple average $\frac{1}{3}(V_0 + V_1 + V_2)$.
\end{itemize}

\textbf{The big picture.} The attention output is not the final embedding --- it is a \emph{context-aware update}. In the full transformer (Document~02), this output is added back to the original embedding via a residual connection: $\mathbf{x} \leftarrow \mathbf{x} + \mathrm{Attention}(\mathbf{x})$. So each position's embedding starts as ``what this token means in isolation'' and gets enriched with ``what the surrounding context says is relevant.'' Stack multiple layers of this, and the embeddings progressively absorb more and more contextual information.

\tip{Since $d_k$ equals the embedding dimension, the attention output has the same shape as the input --- $(T, d_k)$ --- so the residual add works directly.}

\bigskip\hrule\bigskip

\section{Multi-Head Attention}

A single set of $Q, K, V$ projections computes a single attention pattern: one $(T, T)$ matrix of weights. But ``relevance'' is not one-dimensional. Position $j$ might be relevant to position $i$ for syntactic reasons (it is the verb's subject), for semantic reasons (it discusses the same topic), for positional reasons (it is the previous word), and so on.

Multi-head attention addresses this by running multiple attention computations in parallel, each with its own learned projections, and then combining the results.

\textbf{The math.} We now write $d_{\text{model}}$ for the full embedding dimension (what we called $d$ in Sections~2--6). Given input $X$ of shape $(T, d_{\text{model}})$, and $h$ heads:

For each head $i$ (where $i = 1, \ldots, h$):
\[
\text{head}_i = \mathrm{Attention}(X W_i^Q,\; X W_i^K,\; X W_i^V)
\]
where $W_i^Q, W_i^K, W_i^V$ are each $(d_{\text{model}}, d_k)$ projection matrices and $d_k = d_{\text{model}} / h$. Note that $d_k$ now means the \emph{per-head} dimension, smaller than the full embedding dimension we called $d_k$ in Sections~3--6.

Then concatenate and project:
\[
\mathrm{MultiHead}(X) = \mathrm{Concat}(\text{head}_1, \ldots, \text{head}_h)\, W^O
\]
where $W^O$ is $(d_{\text{model}}, d_{\text{model}})$.

Each head operates in a $d_k$-dimensional subspace. Head 1 might learn to attend based on syntactic structure; head 2 might attend to nearby positions; head 3 might track coreference. The output projection $W^O$ learns to combine these different views.

The following snippets show how multi-head attention works step by step inside \code{CausalSelfAttention} in \code{src/model.py}, which we will create in Document~04.

\begin{enumerate}
  \item \textbf{Project to Q, K, V.} Three separate weight matrices $W^Q, W^K, W^V$, each of shape $(d_{\text{model}}, d_{\text{model}})$, produce the queries, keys, and values:

\concepttag
\begin{lstlisting}
q = self.W_Q(x)  # (B, T, d_model)
k = self.W_K(x)  # (B, T, d_model)
v = self.W_V(x)  # (B, T, d_model)
\end{lstlisting}

  \item \textbf{Reshape into heads.} We reshape each $(B, T, d_{\text{model}})$ tensor into $(B, T, n_{\text{head}}, d_k)$ and then transpose to $(B, n_{\text{head}}, T, d_k)$:

\begin{lstlisting}
q = q.view(B, T, self.n_head, self.head_dim).transpose(1, 2)
k = k.view(B, T, self.n_head, self.head_dim).transpose(1, 2)
v = v.view(B, T, self.n_head, self.head_dim).transpose(1, 2)
\end{lstlisting}

Now the tensor has shape $(B, n_{\text{head}}, T, d_k)$. The $n_{\text{head}}$ dimension acts as an additional batch dimension --- the $(T, d_k)$ attention computation runs independently for each head.

  \item \textbf{Run attention.} Standard scaled dot-product attention over the last two dimensions. Each head computes its own $(T, T)$ weight matrix and its own $(T, d_k)$ output:

\begin{lstlisting}
scores = torch.matmul(q, k.transpose(-2, -1)) / math.sqrt(head_dim)
weights = F.softmax(scores, dim=-1)
y = torch.matmul(weights, v)
\end{lstlisting}

This is the same formula from Section~6, now applied per head. Section~8 adds the causal mask that restricts each position to only attend to the past.

  \item \textbf{Reassemble.} After attention, \code{y} has shape $(B, n_{\text{head}}, T, d_k)$. We transpose to $(B, T, n_{\text{head}}, d_k)$ and then merge the last two dimensions back into $d_{\text{model}}$ (since $n_{\text{head}} \times d_k = d_{\text{model}}$):

\begin{lstlisting}
y = y.transpose(1, 2).contiguous().view(B, T, d_model)
\end{lstlisting}

\tip{\code{.contiguous()} is a PyTorch bookkeeping call, not a mathematical operation. \code{.transpose()} does not rearrange data in memory --- it only changes how indices map to the underlying storage. But \code{.view()} requires the data to be laid out contiguously (no gaps or out-of-order strides). \code{.contiguous()} copies the tensor into a fresh, contiguous block so that \code{.view()} can safely reinterpret the shape.}

  \item \textbf{Output projection.} A final linear layer mixes the concatenated head outputs:

\begin{lstlisting}
y = self.c_proj(y)  # (B, T, d_model)
\end{lstlisting}
\end{enumerate}

\textbf{Parameter count.} We have four weight matrices of shape $(d_{\text{model}}, d_{\text{model}})$: $W^Q$, $W^K$, $W^V$, and the output projection $W^O$. Total: $4 \cdot d_{\text{model}}^2$ parameters (ignoring biases). For our BabyGPT with $d_{\text{model}} = 384$, that is $4 \times 384^2 = 589{,}824$ parameters per attention layer.

You can see this exact implementation in \code{src/model.py} in the \code{CausalSelfAttention} class. The three projections are \code{self.W\_Q}, \code{self.W\_K}, \code{self.W\_V}, and the output projection is \code{self.c\_proj}. The attention computation uses the same scaling, masking, softmax, and matrix multiply from Section~6.

\bigskip\hrule\bigskip

\section{Causal (Autoregressive) Masking}

Everything so far describes \textbf{bidirectional} attention: every position can attend to every other position. This is fine for tasks where you have the full input (e.g., classification, encoding). But for language modeling --- predicting the next token given the previous ones --- position $t$ must not see positions $t+1, t+2, \ldots, T$. That would be cheating.

The solution is a \textbf{causal mask}: a lower-triangular matrix of ones.

Causal mask pattern:

\outputtag
\begin{lstlisting}[language={}]
1 0 0 0
1 1 0 0
1 1 1 0
1 1 1 1
\end{lstlisting}

Row $t$ has ones in columns 0 through $t$ and zeros elsewhere. Position 0 can only see itself. Position 1 can see positions 0 and 1. Position $T-1$ can see everything.

We apply this mask to the attention scores \textbf{before} softmax. Where the mask is 0 (future positions), we set the score to $-\infty$:

\demotag
\begin{lstlisting}
mask = torch.tril(torch.ones(T, T))  # lower-triangular
scores = scores.masked_fill(mask == 0, float('-inf'))
\end{lstlisting}

\textbf{Why $-\infty$ and not 0?} Because of how softmax works. If we set a score to 0, then $\exp(0) = 1$, which is a perfectly valid nonzero weight --- the model would still attend to the future position. We need the weight to be exactly zero, and the only way to get softmax to output zero is to feed it $-\infty$: $\exp(-\infty) = 0$. The masked position contributes nothing to the weighted average.

After masking and softmax, the weight matrix looks like:

\outputtag
\begin{lstlisting}[language={}]
[1.00  0.00  0.00  0.00]
[0.45  0.55  0.00  0.00]
[0.20  0.30  0.50  0.00]
[0.10  0.25  0.35  0.30]
\end{lstlisting}

Each row sums to 1, and all entries above the diagonal are exactly 0. Every position attends only to itself and to the past.

In \code{src/model.py}, the causal mask is constructed with \code{torch.tril} and applied before softmax, exactly as shown above:

This is how it appears in the final \code{src/model.py}.

\concepttag
\begin{lstlisting}
scores = torch.matmul(q, k.transpose(-2, -1)) / math.sqrt(self.head_dim)
mask = torch.tril(torch.ones(T, T, device=x.device, dtype=torch.bool))
scores = scores.masked_fill(~mask, float('-inf'))
weights = F.softmax(scores, dim=-1)
weights = self.attn_dropout(weights)
y = torch.matmul(weights, v)
\end{lstlisting}

This is the complete attention computation from \code{CausalSelfAttention.forward} --- scaling, masking, softmax, dropout, and the weighted sum of values, all visible.

\bigskip\hrule\bigskip

\section{Looking Ahead}

We have now covered the core computational primitive of the transformer: scaled dot-product attention, extended to multiple heads, with causal masking for autoregressive modeling.

But attention alone is not a transformer. It is one sublayer in a larger block that also includes:

\begin{itemize}
  \item \textbf{Residual connections} --- the input to each sublayer is added to its output, enabling gradient flow through deep networks.
  \item \textbf{Layer normalization} --- stabilizes training by normalizing activations.
  \item \textbf{A feed-forward network} --- a position-wise MLP that gives each position capacity to process its attended representation.
\end{itemize}

The next document, \textit{02~---~The Transformer Block}, shows how these components wrap around attention to form the repeating unit of the transformer architecture.

\end{document}
